\section{Primitive squares in two odd extensions}
This section reproduces the project's earlier accepted primitive-square
proof, independently audited in phase267. Its new placement in the
common theory is not a new answer to a source question. The retained
analytic input is the multiplicative-character bound in
Cohen--Kapetanakis--Reis, \emph{The existence of $F_q$-primitive points
on curves using freeness}, Theorem~5 (printed p.~643), together with
the cyclic-group primitive-element indicator. The exact primary and
prior proof bindings accompany the scope map. The following particular
norm construction and finite reduction are the previously proved work.

\begin{theorem}\label{E.PrimitiveSquare}
Let $p$ be an odd prime and $e_1,e_2>1$ odd integers. There are primitive
$j_i\in F_{p^{e_i}}$ such that, with $u=j_1+j_1^{-1}$ and
$v=j_2+j_2^{-1}$,
\[
 u^2v^2-4(u^2+v^2)
\]
is a nonzero square in $F_{p^{\operatorname{lcm}(e_1,e_2)}}$.
The degrees may be equal or nested.
\end{theorem}
\begin{proof}
Reorder the degrees and put $g=\gcd(e_1,e_2)$, $q_0=p^g$,
$a=e_1/g\le b=e_2/g$. Then $a,b$ are coprime and odd. Write
$E=F_{q_0^a}$, $K=F_{q_0^b}$, $L=F_{q_0^{ab}}$.
Fix any primitive $j_1\in E$ and set $s=(j_1+j_1^{-1})^2$.
One has $s\ne0,4$, since otherwise $j_1$ has order at most four.
Also $s$ has degree $a$ over $F_{q_0}$: if it lies in a proper
intermediate field, the equations $u^2=s$ and $j_1^2-uj_1+1=0$
put first $u$ and then $j_1$ in that same field, because the remaining
extension degree is odd. This contradicts primitiveness.

Define
\[
 f_s(X)=(s-4)(X^2+1)^2-4sX^2,
 \qquad P(X)=\prod_{i=0}^{a-1}f_{s^{q_0^i}}(X).
\]
For $x\ne0$ the desired quantity with $j_2=x$ is $f_s(x)/x^2$.
The quadratic in $X^2$ defining $f_s$ has nonzero constant and leading
coefficients $s-4$ and discriminant $64s\ne0$. Thus $f_s$ has four
distinct nonzero roots. Distinct quartics have disjoint roots: subtracting
gives $(s-t)(X^2-1)^2$, while $f_s(\pm1)=-16\ne0$.
Consequently $P\in F_{q_0}[X]$ is squarefree of degree $D=4a$ and
$P(0)\ne0$. Since $a,b$ are coprime, the $q_0^b$-Frobenius permutes
all these $a$ conjugates. For $x\in K$ this proves
\[
 N_{L/K}(f_s(x))=P(x),\qquad
 \chi_L(f_s(x)/x^2)=\chi_K(P(x)),
\]
where $\chi$ is the quadratic character extended by zero.
If $a>1$, $P$ has no root in $K$: a root would give
$s^{q_0^i}=4(x^2+1)^2/(x^2-1)^2\in E\cap K=F_{q_0}$,
contradicting its degree. When $a=1$, at most four roots remain.

Put $Q=|K|$, $n=Q-1$, $W(n)=2^{\omega(n)}$ and
$\vartheta(n)=\varphi(n)/n$. For every multiplicative character $\psi$,
\[
 \left|\sum_{x\in K^\times}\psi(x)\chi_K(P(x))\right|
 \le D\sqrt Q.                                             \tag{20}
\]
Indeed, for nontrivial $\psi=\eta^k$, $1\le k<n$, with $\eta$ a
generator of the character group, the summand is
$\eta(x^kP(x)^{n/2})$. Its polynomial has $D+1$ distinct roots and
at a root of $P$ has multiplicity exactly $n/2$, so is not a scalar
times an $n$th power. The retained bound gives $D\sqrt Q$.
For trivial $\psi$, apply the quadratic-character bound on all of $K$
and remove zero, giving $(D-1)\sqrt Q+1\le D\sqrt Q$.
All characters, including the trivial one, are extended by zero at zero.

The primitive-element indicator is
\[
 {\bf1}_{\rm primitive}(x)=\vartheta(n)
 \sum_{d\mid n}\frac{\mu(d)}{\varphi(d)}
       \sum_{\operatorname{ord}\psi=d}\psi(x).
\]
Character orthogonality in the cyclic group $K^\times$ proves this
formula. There are $\varphi(d)$ characters of order $d$ and $W(n)$
squarefree divisors. If $Z$ counts primitive roots of $P$, (20) therefore
makes the number of primitive $x$ with $P(x)$ a nonzero square positive
provided
\[
 Q-1>D W(Q-1)\sqrt Q+Z/\vartheta(Q-1).
\]
For $a>1$ it suffices that
\[
 Q-1>4aW(Q-1)\sqrt Q;                                    \tag{21}
\]
for $a=1$, since $Z\le4$ and $1/\vartheta(n)\le W(n)$, it suffices that
\[
 Q-1>4W(Q-1)(\sqrt Q+1).                                 \tag{22}
\]

We give the full finite reduction. For every positive integer $n$,
\[
 W(n)\le \frac{64}{30030^{1/4}}n^{1/4}<5n^{1/4}.
\]
Each prime at least $17$ contributes $2\le\ell^{1/4}$; only
$2,3,5,7,11,13$, whose product is $30030$, can enlarge the ratio.
Repeated prime powers decrease it, and $64^4<5^4\cdot30030$.
Thus (21) holds when $Q>(21a)^4$.
In the nonnested case $3\le a\le b-2$, $b\ge5$. For odd $b\ge23$,
$3^b>[21(b-2)]^4$, first by direct comparison at $23$ and then because
the ratio increases on replacing $b$ by $b+2$ by a factor at least
$9(21/23)^4>1$. For odd $5\le b\le21$, it remains only to examine
$q_0^b\le[21(b-2)]^4$ with $q_0=p^g$, $g$ odd.
The possible bases are at most $27,13,9,5,5,3,3,3,3$ in this order;
the only composite base of odd prime-power exponent is $27$.
The exact integer comparison
\[
 (Q-1)^2>16a^2W(Q-1)^2Q
\]
for all $66$ remaining parameter triples leaves precisely
\[
 (q_0,a,b)=(3,3,5),(5,3,5),(3,3,7),(3,5,7),(3,5,9),(3,7,9).
\]
This exhaustive bounded arithmetic is part of the required program,
using exact trial division and integer comparisons.

For $a=1$, (22) follows from $Q>21^4$: writing $t=Q^{1/4}>21$,
$(t-20)t^3>t^3>20t+1$ gives $Q-1>20t^3+20t$.
There are exactly $23$ odd prime-power fields $p^e\le21^4$, $e\ge3$
odd. The exact tests $Q-1-4W>0$ and
$(Q-1-4W)^2>16W^2Q$ leave only $Q=3^3,3^5,5^3,7^3,11^3$.
Their exponents are prime, so the divisor $e_1>1$ of $e_2$ equals $e_2$.

Finally put $z_e=Z(p^e)$ in the bound Conway-field representation,
choose $j_1=z_{e_1}$ and $j_2=z_{e_2}^k$, and use the following complete
exceptional list. The displayed exponent is even and nonzero modulo
$p^{e_2}-1$, and $k$ is coprime to that integer.
{\renewcommand{\arraystretch}{1.4}
\[
\begin{array}{ccc|r|r}
 p&e_1&e_2&k& P(j_2)\text{ or }f_s(j_2)\\\hline
 3&3&5&7&z_5^{228}\\5&3&5&1&z_5^{2802}\\
 3&3&7&1&z_7^{1336}\\3&5&7&1&z_7^{406}\\
 3&5&9&5&z_9^{18778}\\3&7&9&1&z_9^{17426}\\
 3&3&3&1&z_3^6\\3&5&5&3&z_5^{200}\\
 5&3&3&1&z_3^{84}\\7&3&3&1&z_3^6\\11&3&3&3&z_3^{906}
\end{array}
\]}
The complete polynomial coefficients and defining irreducible polynomials
are carried in \path{psl2_square_witnesses.json}. The required
\path{verify_square_witnesses.py} checks the finite-field arithmetic
in explicit polynomial quotient rings, independently of GAP; the
required \path{psl2_square_bound_check.py} checks both bounded
reductions. The old GAP witness generators are retained separately as
validation. These eleven exact certificates close all exceptions to
the uniform inequality and hence all parameters of the theorem.
\end{proof}
